\chapter*{О ГОСе\footnotemark}
\addcontentsline{toc}{chapter}{О ГОСе}

\footnotetext{Описаны реальные события 2026 года. Возможно, что-то поменялось, и информация здесь уже не актуальна.}
В первую очередь, нужно рассказать о процедуре проведения итогового экзамена по математике в МФТИ глазами студента третьего курса. 

Процедура проведения этого экзамена почти ничем не отличается от обычных экзаменов кафедры высшей математики. 
Успех на этом экзамене зависит, прежде всего, от вашей удачи, затем от уровня подготовки и, напоследок, от того, насколько вы умеете <<вертеться>>, впрочем, как и на протяжении всей жизни на Физтехе...

Сначала Вас ждёт письменная работа, в которой будет много задач, порядка 18. На нее отводится три астрономических часа, т.е.~180 минут

Конечно же, рекомендуется тщательно подготовиться к письменному экзамену, прорешать множество вариантов, посмотреть консультации преподавателей по решению задач и отнестись к ним лишь рекомендательно, потому что каждый год кафедра высшей математики преподносит некоторые сюрпризы, например, появляются задачи, о которых на консультациях преподаватели говорят: <<Не будет>>, <<Маловероятно>>, <<Это слишком сложно для ГОСа>>. Получить достойные баллы при должной подготовке вполне реально. Разбалловку рассчитывают, основываясь на результатах всех студентов (подгоняя под распределение Гаусса), в наш год было так: средний балл около 15 баллов из 68. Контрольная оказалась довольно сложно и коэффициент перевода был 1:1.

С этого года оценка за экзамен определяется шкалой БРС, уже привычной Вам по предыдущим экзаменам от КВМ. Но с одним отличием: 40 из 120 баллов можно получить за работу в семестрах (средний балл за устные экзамены от КВМ). За писбменную работу - всего-лишь 30 очков, и остальные 50 за устный ответ.

Перейдем к обсуждению устного экзамена. Многие после ГОСа по физике удивляются, что устный экзамен по математике проходит в больших аудиториях, таких как Актовый Зал, Большая Физическая, 117 ГК и прочие, но на самом деле, ничего удивительного. Как оказалось, это самый обычный экзамен от кафедры вышмата с некоторыми особенностями, о которых ниже. Конечно же, этот экзамен "--- это один из самых больших по объему материалов для подготовки, поэтому усердно работайте, постарайтесь хорошо подготовиться, надеюсь, это пособие вам поможет.

С этого года пересдачи на ГОСе "--- не такая и редкость, но всё зависит от комисии: на некоторых факультетах было 5-7 пересдач, что всё равно очень мало, но раньше их практически не ыбло и получить неуд за ГОС было крайне сложно.
На ГОСе нельзя проситься к преподавателям, а они бывают разные: у всех разное отношение к студентам, к самому итоговому экзамену, особенно полсе того, как его сделали совершенно обычным экзаменом. В коцне вам просто назовут оценку и посчитают итоговый БРС: без рбщего совещания комиссий, кофебреков и больших перерывов.

Так и заканчивается учебный год. Целый период жизни на Физтехе. Все прощаются с кафедрой высшей математики. Все уходят на летний отдых или по делам.

\mbox{}

\noindent\textit{Диденко Андрей, Лобачев Александр}


